\documentclass[10pt, letterpaper]{article}

\usepackage[margin=1in]{geometry}
\usepackage[T1]{fontenc}
\usepackage{lmodern}
\usepackage{amsmath, amssymb, amsthm}
\usepackage{booktabs}
\usepackage{graphicx}
\usepackage{xcolor}
\usepackage{hyperref}
\usepackage{enumitem}
\usepackage{caption, subcaption}
\definecolor{groupmoe}{RGB}{31,119,180}
\definecolor{stdmoe}{RGB}{255,127,14}
\definecolor{baseline}{RGB}{44,160,44}

\newcommand{\R}{\mathbb{R}}
\newcommand{\Rg}{R(g)}
\newcommand{\gmoe}{\textsc{Group-MoE}}

\title{Group-MoE: Learned Dispatch to Algebraic Fixed-Function Units}
\author{Robb Walters}
\date{April 2026 \quad\textbar\quad arXiv Preprint}

\begin{document}
\maketitle

\begin{abstract}
Neural networks that encounter symmetry either enforce it rigidly in every layer or ignore it entirely. We propose \gmoe{}, a mixture-of-experts architecture where each expert implements a finite group representation in the irreducible representation (irrep) basis, and a learned router detects which symmetry applies. This gives models \emph{algebraic fixed-function units}---analogous to dedicated hardware accelerators---that the network learns to invoke selectively. We introduce the \emph{complement split}, a controlled evaluation methodology that isolates symmetry transfer from memorization, and validate the architecture on synthetic tasks. Our experiments demonstrate: (1)~complement transfer from seen to unseen orderings (+16 percentage points for $S_2$, +3.9pp for $S_3$ beyond a matched standard MoE), (2)~near-perfect zero-shot compositional generalization from transpositions to 3-cycles via irrep composition, (3)~correct router discrimination between non-nested groups, and (4)~zero-degradation compatibility as a transformer FFN replacement. We analyze when algebraic structure outperforms lookup-table memorization and identify the combinatorial scaling regime where group experts should provide their greatest advantage.
\end{abstract}

%----------------------------------------------------------------------
\section{Introduction}
\label{sec:intro}

Symmetry is ubiquitous in the data that neural networks process. Images have rotational and reflective symmetries. Molecules have point-group symmetries. Language has permutation symmetries: ``Alice likes Bob'' and ``Bob is liked by Alice'' describe the same relationship regardless of entity ordering. The question is not whether symmetry exists, but how a model should handle it.

The field has converged on two approaches, each with significant limitations. \emph{Geometric deep learning}~\cite{bronstein2021geometric, cohen2016group} builds equivariance into every layer, ensuring that symmetry transformations of the input produce corresponding transformations of the output. This is correct by construction but rigid: the model cannot mix symmetric and non-symmetric computation, and applying the wrong symmetry is as harmful as applying none. \emph{Standard architectures} (MLPs, transformers) ignore symmetry entirely and hope the right behavior emerges from training data. Our prior work~\cite{walters2025ordering} showed that large language models do achieve \emph{functional} equivariance---correct outputs under input permutation---but without developing \emph{structural} equivariance---group representations in their latent space. They memorize the multiplication table rather than learning algebra.

This paper proposes a third path: \textbf{group representations as optional expert modules} with learned dispatch. The analogy to hardware design is precise. Just as processors evolved from general-purpose CPUs to include fixed-function accelerators (FFT units, matrix multiply engines, cryptographic co-processors), neural architectures can include \emph{algebraic fixed-function units}---modules hardwired in the irreducible representation basis of a finite group---that are activated only when the router detects the corresponding symmetry. The irrep basis provides the same kind of structural compression that fixed-function hardware provides over general-purpose computation: a $k \times k$ block-diagonal transform instead of a $d \times d$ dense matrix, with exact composition guaranteed by the group homomorphism property.

Our contributions:
\begin{enumerate}[leftmargin=*]
    \item The \gmoe{} architecture: group representations in the irrep basis as expert modules, combined with a learned symmetry router and residual blending (Section~\ref{sec:arch}).
    \item The \emph{complement split}: a controlled evaluation methodology that isolates symmetry transfer from memorization by ensuring every test example requires generalization from a seen ordering to an unseen one (Section~\ref{sec:complement}).
    \item Experimental validation on five synthetic tasks, each answering one question about the architecture (Section~\ref{sec:experiments}).
    \item Analysis of when algebraic structure outperforms lookup-table memorization, connecting our findings to the broader question of scaling (Section~\ref{sec:analysis}).
\end{enumerate}

%----------------------------------------------------------------------
\section{Architecture}
\label{sec:arch}

\subsection{Group Expert Module}

The core building block of \gmoe{} is the \emph{group expert}, which applies a group representation in a learned subspace of the activation space. For a finite group $G$ with irreducible representations of total dimension $k$, the expert implements:
\begin{equation}
    x \xrightarrow{P} z \xrightarrow{\Rg} z' \xrightarrow{P^\dagger} x'
    \label{eq:expert}
\end{equation}
where $x \in \R^d$ is the input activation, $P \in \R^{k \times d}$ is a learned projection into the symmetry-active subspace, $\Rg \in \R^{k \times k}$ is the irrep matrix for group element $g$ (fixed, not learned), and $P^\dagger \in \R^{d \times k}$ is the injection back to the full space.

The irrep matrix $\Rg$ is block-diagonal by construction:
\begin{equation}
    \Rg = \begin{pmatrix} \rho_1(g) & & \\ & \rho_2(g) & \\ & & \ddots \end{pmatrix}
\end{equation}
where each $\rho_i$ is an irreducible representation. For $S_3$, the symmetric group on three elements, the irreps are: trivial (1D), sign (1D), and standard (2D), giving $k = 4$. For a model with $d = 128$, this provides $32\times$ compression relative to a full $d \times d$ matrix.

The critical property is \textbf{exact composition}: $R(g_1) R(g_2) = R(g_1 g_2)$ by the group homomorphism. This means the expert's behavior on composed group elements is determined entirely by its behavior on generators---a property no learned transform can guarantee.

\subsection{Symmetry Router}

The router is a lightweight MLP that maps each activation to a distribution over routing options: one pass-through option plus one option per group element across all available groups. For a model with groups $G_1, \ldots, G_m$:
\begin{equation}
    \text{options} = \{{\text{pass-through}}\} \cup \bigcup_{i=1}^{m} \{g \mid g \in G_i\}
\end{equation}

The router uses hard routing (argmax) with soft confidence blending:
\begin{equation}
    x' = x + \alpha \cdot (\text{expert}(x, g) - x)
\end{equation}
where $\alpha$ is the router's confidence score for the selected option. Pass-through tokens ($\alpha = 0$ or pass-through selected) pass unchanged, ensuring that symmetry is only applied when detected.

\subsection{Load Balancing}

Following the Switch Transformer~\cite{fedus2022switch}, we add an auxiliary load-balancing loss to prevent expert collapse:
\begin{equation}
    \mathcal{L}_{\text{balance}} = N \sum_{i=1}^{N} f_i \cdot p_i
\end{equation}
where $N$ is the number of routing options, $f_i$ is the fraction of tokens assigned to option $i$ (hard), and $p_i$ is the mean routing probability for option $i$ (soft). This encourages uniform utilization of all options without forcing specific routing patterns.

\subsection{Integration}

The \gmoe{} layer is a residual module that handles both $(B, d)$ and $(B, L, d)$ inputs by flattening the sequence dimension internally. This makes it a drop-in replacement for any feedforward block in a transformer or MLP architecture.

Table~\ref{tab:groups} lists the implemented groups with their irrep dimensions and parameter counts.

\begin{table}[h]
\centering
\caption{Implemented group representations.}
\label{tab:groups}
\begin{tabular}{llccc}
\toprule
Group & Order & Irreps & $k$ & Expert params ($d{=}128$) \\
\midrule
$\mathbb{Z}_2$ & 2 & trivial + sign & 2 & 512 \\
$\mathbb{Z}_3$ & 3 & trivial + standard(2D) & 3 & 768 \\
$S_2$ & 2 & trivial + sign & 2 & 512 \\
$S_3$ & 6 & trivial + sign + standard(2D) & 4 & 1,024 \\
\bottomrule
\end{tabular}
\end{table}

%----------------------------------------------------------------------
\section{The Complement Split}
\label{sec:complement}

Standard random train/test splits do not isolate symmetry-specific generalization. A model might generalize from $(a, +, b)$ to $(b, +, a)$ through embedding interpolation rather than by learning commutativity. To measure the group expert's structural contribution, we need a split where every test example \emph{requires} symmetry-specific transfer.

\textbf{Definition.} For a symmetric function $f$ on $n$ inputs with symmetry group $G$, the \emph{complement split} assigns exactly one ordering per $G$-orbit to training and all remaining orderings to test. Formally, for each orbit $\mathcal{O} = \{g \cdot (a_1, \ldots, a_n) \mid g \in G\}$, we select one representative $\sigma \in \mathcal{O}$ for training; the remaining $|\mathcal{O}| - 1$ elements go to test.

\textbf{For $S_2$ on pairs $\{a, b\}$:} train on $(a, +, b)$, test on $(b, +, a)$. Every test example is a reversal of a training example.

\textbf{For $S_3$ on triples $\{a, b, c\}$:} train on 1 ordering, test on the other 5. The generalization ratio is 1:5---substantially harder than $S_2$'s 1:1.

\textbf{Composition split} (variant): train on identity + transpositions (4 orderings), test on 3-cycles (2 orderings). This tests compositional generalization: the 3-cycle orderings are compositions of transpositions, and the irrep basis computes them by construction ($R(\sigma_1 \sigma_2) = R(\sigma_1) R(\sigma_2)$), while a model without group structure must learn them independently.

In all cases, a non-symmetric control operation is split randomly, providing a baseline for the model's general learning capacity on the same data.

%----------------------------------------------------------------------
\section{Experiments}
\label{sec:experiments}

We present five experiments, each answering one question. All use the same base architecture: opaque embeddings (one per number, shared across positions), a concatenation layer, and residual MLP blocks, followed by the group expert layer and a scalar regression head. Unless noted, $d = 128$, 2 MLP blocks, AdamW with learning rate $3 \times 10^{-4}$, weight decay $10^{-2}$, balance loss $\alpha = 0.01$.

\subsection{Does complement transfer work? ($S_2$, Arithmetic)}
\label{sec:exp-s2}

\textbf{Task.} $(a, \text{op}, b) \to \text{result}$, op $\in \{+, -\}$, $a, b \in [0, 20)$. Complement split on addition pairs: train on $(a, +, b)$, test on $(b, +, a)$. Subtraction split randomly.

\textbf{Result.} \gmoe{} achieves 48.4\% complement accuracy vs.\ 31.6\% for the baseline---a 53\% relative improvement. The effect is consistent across 3 random seeds (Table~\ref{tab:s2-seeds}).

\begin{table}[h]
\centering
\caption{$S_2$ complement transfer across seeds (num\_range=20).}
\label{tab:s2-seeds}
\begin{tabular}{lccc}
\toprule
Seed & \gmoe{} & Baseline & Relative gain \\
\midrule
42 & \textbf{48.4\%} & 31.6\% & +53\% \\
123 & \textbf{45.8\%} & 27.9\% & +64\% \\
7 & \textbf{35.3\%} & 33.2\% & +6\% \\
\bottomrule
\end{tabular}
\end{table}

\textbf{Router behavior.} The router routes 95\% of addition examples to the $S_2$ expert and 35\% of subtraction examples---clean operation discrimination. Per-pair analysis shows \gmoe{} exclusively solves 2.5$\times$ more pairs than the baseline (40 vs.\ 16 out of 190).

\textbf{Balance loss ablation.} $\alpha = 0.01$ is the sweet spot. $\alpha = 0$ still produces complement transfer (+30\% relative) but with weaker router discrimination. $\alpha = 0.1$ forces excessive routing and hurts performance.

% TODO: figure needed — S_2 complement heatmaps (error, advantage, routing)

\subsection{Is it the group structure or just the routing? ($S_3$, Three-Way)}
\label{sec:exp-threeway}

A skeptic could attribute the advantage to the routing architecture rather than the group structure. We add a controlled middle ground: \textsc{StandardMoE}, which has the same router, the same project-inject dimensions ($d \to k \to d$), the same residual blending---but replaces the fixed irrep matrix $\Rg$ with a \emph{learned} $k \times k$ matrix $W$. The three models have nearly identical parameter counts ($\sim$337K).

\textbf{Task.} $(a, b, c, \text{op}) \to \text{result}$, op $\in \{\text{sum}, \text{nonsym}\}$, num\_range=15. Two split modes: complement (1 ordering trains, 5 test) and composition (identity + transpositions train, 3-cycles test).

\begin{table}[h]
\centering
\caption{Three-way comparison on $S_3$ ternary task (num\_range=15).}
\label{tab:threeway}
\begin{tabular}{lcc}
\toprule
Model & Complement (1$\to$5) & Composition (4$\to$2) \\
\midrule
\gmoe{} (irrep $\Rg$) & \textbf{92.1\%} & \textbf{99.0\%} \\
StandardMoE (learned $W$) & 88.2\% & 98.1\% \\
Baseline (no expert) & 86.1\% & 96.2\% \\
\bottomrule
\end{tabular}
\end{table}

\textbf{Decomposition.} On the complement split, the group structure contributes +3.9pp beyond what the routing architecture alone provides (+2.1pp). On the easier composition split, the group contribution is +0.9pp vs.\ +1.9pp from routing. \textbf{The harder the generalization, the more the algebraic structure matters.}

\subsection{Does irrep composition generalize? (Transpositions $\to$ 3-Cycles)}
\label{sec:exp-composition}

The strongest theoretical prediction of group representations: since $R(\sigma_1 \sigma_2) = R(\sigma_1) R(\sigma_2)$, the expert should work for 3-cycles without ever training on them.

\textbf{Setup.} Composition split: train on identity + 3 transpositions per triple (generators of $S_3$), test only on the 2 three-cycle orderings (compositions).

\textbf{Result.} \gmoe{} achieves 98.5--99.0\% on 3-cycle orderings never seen during training. The projection $P$, learned from transposition examples, generalizes because the irrep subspace is shared across all group elements. At num\_range=10, the router achieves perfect discrimination: 100\% $S_3$ routing for the symmetric operation vs.\ 74\% for the non-symmetric control.

\subsection{Can the router discriminate between groups?}
\label{sec:exp-multigroup}

\textbf{Nested groups ($S_2 + S_3$).} With both $S_2$ and $S_3$ experts available, the router routes $\sim$70\% of all operations to $S_3$ regardless of operation type. This is \emph{rational}: since $S_2 \subset S_3$, every $S_2$ transformation is also an $S_3$ transformation. The router discovers the subgroup relationship without supervision.

\textbf{Non-nested groups ($\mathbb{Z}_2 + \mathbb{Z}_3$).} With genuinely disparate groups (neither is a subgroup of the other), the router correctly dispatches $\mathbb{Z}_2$ operations to the $\mathbb{Z}_2$ expert at 76\%. The non-symmetric control routes primarily to pass-through (15\%) or uses the larger expert as general-purpose compute.

This confirms the ASIC analogy: the router learns which algebraic unit to invoke when the units have non-overlapping capabilities.

\subsection{Does it work in a transformer?}
\label{sec:exp-transformer}

\textbf{Setup.} 4-layer transformer encoder ($d = 256$, 4 heads). Each of $(a, \text{op}, b, c)$ is a separate token with learned positional embeddings. \gmoe{} replaces the FFN in one block.

\textbf{Result.} All three models (Group-MoE, StandardMoE, Baseline) reach 100\% complement accuracy. Self-attention handles permutation mixing natively for short sequences, providing functional equivariance without group structure. \gmoe{} is a zero-degradation drop-in FFN replacement, confirming architectural compatibility. The group expert's marginal value will manifest on tasks where attention alone cannot resolve the symmetry.

%----------------------------------------------------------------------
\section{When Does Algebraic Structure Beat Memorization?}
\label{sec:analysis}

Our experiments reveal a consistent pattern that we frame as the \textbf{lookup table (LUT) vs.\ algebra} tradeoff.

\textbf{Memorization wins when:} (1)~the function decomposes into per-element contributions (e.g., sum = per-number values), so the effective table size is $O(n)$ not $O(n^k)$; (2)~the model has enough capacity relative to the table size; (3)~attention provides permutation mixing natively; (4)~enough orderings are observed that interpolation suffices.

\textbf{Algebraic structure helps when:} (1)~the complement split isolates symmetry transfer from memorization; (2)~the model's backbone does not already handle permutation mixing; (3)~the effective lookup table has gaps that the group structure can fill; (4)~the irrep composition property enables zero-shot generalization to unseen composed elements.

\textbf{The scaling prediction.} The group $S_n$ has $n!$ elements. For $S_{10}$, this is 3.6 million; for $S_{20}$, $2.4 \times 10^{18}$. No model memorizes $10^{18}$ entries, but the irrep basis for $S_{20}$ remains structured and compressible. Our synthetic experiments operate in the small-table regime where memorization is competitive. The compelling case for \gmoe{} is the \emph{large-table regime}---large groups, long sequences, combinatorial state spaces---where the factorial growth of the lookup table defeats memorization.

The developmental analogy is apt: humans learn multiplication tables by rote before understanding algebra. For small tables, rote memorization is strictly easier. Algebra pays off when the table is too large to memorize---and in the limit, it is the \emph{only} strategy that works.

%----------------------------------------------------------------------
\section{Related Work}
\label{sec:related}

\textbf{Geometric deep learning.} Cohen \& Welling~\cite{cohen2016group} introduced group equivariant convolutions; Bronstein et al.~\cite{bronstein2021geometric} unified CNNs, GNNs, and transformers under the equivariance framework; Weiler \& Cesa~\cite{weiler2019general} developed general steerable CNNs using irrep decomposition. All enforce symmetry in every layer. \gmoe{} uses the same irrep machinery but makes it optional via routing.

\textbf{Mixture of experts.} Shazeer et al.~\cite{shazeer2017outrageously} proposed sparse MoE with gating; Fedus et al.~\cite{fedus2022switch} simplified to top-1 routing with load balancing. Our router and balance loss follow this design, but our experts have algebraic structure rather than being generic MLPs. Kang et al.~\cite{kang2025mixture} propose ``Mixture of Group Experts'' using group sparse regularization for expert diversity; despite the shared name, their ``groups'' refer to expert grouping, not group theory.

\textbf{Learned group representations.} MatrixNet~\cite{laird2024matrixnet} learns matrix representations of group elements from data, achieving compositional generalization via learned generator representations. Our approach is complementary: we fix the irrep structure (guaranteeing exact composition) and learn the dispatch (routing).

\textbf{Symmetry detection.} Dehmamy et al.~\cite{dehmamy2021automatic} discover continuous symmetries via Lie algebra convolution; Benton et al.~\cite{benton2020learning} meta-learn invariances. These discover symmetries post-hoc. \gmoe{} provides exact group structure as architectural options with learned selection.

\textbf{Compositional generalization.} Gordon et al.~\cite{gordon2020permutation} hypothesize language compositionality is group-equivariance and build equivariant seq2seq models. We share the hypothesis but provide group structure as optional expert modules rather than mandatory architectural constraints.

\textbf{Functional vs.\ structural equivariance.} Walters~\cite{walters2025ordering} (``Ordering Is Not Invariant'') showed LLMs achieve functional equivariance without structural equivariance. \gmoe{} is the architectural response: providing structural equivariance as an option.

%----------------------------------------------------------------------
\section{Discussion and Future Work}
\label{sec:discussion}

\textbf{The ASIC analogy.} Group experts are cognitive arithmetic units---domain-specific algebraic accelerators hardwired in the irrep basis, dispatched by a learned router. The analogy to hardware evolution is not merely metaphorical: the irrep basis provides the same structural compression that fixed-function hardware provides over general-purpose computation, and the router serves the same dispatch role as a processor's instruction decoder.

\textbf{Scaling to combinatorial regimes.} Our synthetic experiments validate the mechanism but operate where memorization is competitive. The key prediction: for $S_n$ with large $n$, the $n!$ growth of the lookup table will defeat memorization while the irrep basis remains tractable. Testing this requires implementing larger permutation groups or finding natural tasks with large effective symmetry groups.

\textbf{Language modeling.} Natural language contains implicit symmetry: entity permutations, fact reorderings, syntactic transforms. These create effective symmetry groups whose ``multiplication tables'' are too large for rote memorization. Inserting \gmoe{} into language models and testing with entity permutation probes~\cite{walters2025ordering} is the natural next step.

\textbf{Symmetry discovery.} Currently the groups are specified a priori. A compelling extension: learn which groups are useful from data, by treating the irrep structure as differentiable or using a library of candidate groups with learned selection.

\textbf{Approximate symmetry.} Real data has approximate, not exact, symmetry. Soft irrep decomposition with learned tolerance could extend \gmoe{} to noisy settings.

%----------------------------------------------------------------------
\section{Conclusion}
\label{sec:conclusion}

\gmoe{} demonstrates that neural networks can learn to dispatch to algebraic fixed-function units. The complement split methodology cleanly isolates symmetry transfer from memorization. The three-way comparison shows the advantage comes from group structure specifically (+3.9pp on $S_3$ complement), not just routing architecture (+2.1pp). Compositional generalization confirms the theoretical payoff: irrep composition enables near-perfect zero-shot transfer to composed elements (99.0\%). The router learns group-theoretic relationships---subgroup containment, non-nested discrimination---without supervision.

These are proof-of-concept results on synthetic tasks where memorization is competitive. The architecture's value will compound in regimes where the effective symmetry group is too large to memorize---exactly the combinatorial territory where current models struggle. \gmoe{} provides the algebraic tools; identifying the right large-scale application is future work.

%----------------------------------------------------------------------
\begin{thebibliography}{99}

\bibitem{bronstein2021geometric}
M.~M.~Bronstein, J.~Bruna, T.~Cohen, and P.~Veli\v{c}kovi\'{c}.
\newblock Geometric deep learning: Grids, groups, graphs, geodesics, and gauges.
\newblock \emph{arXiv:2104.13478}, 2021.

\bibitem{cohen2016group}
T.~S.~Cohen and M.~Welling.
\newblock Group equivariant convolutional networks.
\newblock In \emph{ICML}, 2016.

\bibitem{weiler2019general}
M.~Weiler and G.~Cesa.
\newblock General {E(2)}-equivariant steerable {CNNs}.
\newblock In \emph{NeurIPS}, 2019.

\bibitem{shazeer2017outrageously}
N.~Shazeer, A.~Mirhoseini, K.~Maziarz, A.~Davis, Q.~Le, G.~Hinton, and J.~Dean.
\newblock Outrageously large neural networks: The sparsely-gated mixture-of-experts layer.
\newblock In \emph{ICLR}, 2017.

\bibitem{fedus2022switch}
W.~Fedus, B.~Zoph, and N.~Shazeer.
\newblock Switch {T}ransformers: Scaling to trillion parameter models with simple and efficient sparsity.
\newblock \emph{JMLR}, 23(120):1--40, 2022.

\bibitem{kang2025mixture}
L.~Kang, J.~Li, M.~Tian, and H.~Huang.
\newblock Mixture of group experts for learning invariant representations.
\newblock \emph{arXiv:2504.09265}, 2025.

\bibitem{laird2024matrixnet}
L.~Laird, C.~Hsu, A.~Bapat, and R.~Walters.
\newblock {MatrixNet}: Learning over symmetry groups using learned group representations.
\newblock In \emph{NeurIPS}, 2024.

\bibitem{dehmamy2021automatic}
N.~Dehmamy, R.~Walters, Y.~Liu, D.~Wang, and R.~Yu.
\newblock Automatic symmetry discovery with {L}ie algebra convolutional network.
\newblock In \emph{NeurIPS}, 2021.

\bibitem{benton2020learning}
G.~Benton, M.~Finzi, P.~Izmailov, and A.~G.~Wilson.
\newblock Learning invariances in neural networks from training data.
\newblock In \emph{NeurIPS}, 2020.

\bibitem{gordon2020permutation}
J.~Gordon, D.~Lopez-Paz, M.~Baroni, and D.~Bouchacourt.
\newblock Permutation equivariant models for compositional generalization in language.
\newblock In \emph{ICLR}, 2020.

\bibitem{walters2025ordering}
R.~Walters.
\newblock Ordering is not invariant: Probing structural equivariance in large language models.
\newblock 2025.

\end{thebibliography}

\appendix

\section{Group Representation Details}
\label{app:groups}

All irrep matrices are verified to satisfy the group composition property: $R(g_1)R(g_2) = R(g_1 g_2)$ for all pairs $(g_1, g_2)$. Full multiplication tables and irrep matrices are available in the codebase (\texttt{src/groups/representations.py}).

\section{Hyperparameter Sensitivity}

\textbf{Balance loss $\alpha$.} Tested at $\alpha \in \{0, 0.01, 0.1\}$ on $S_2$ arithmetic complement split. $\alpha = 0.01$ is the sweet spot: $\alpha = 0$ provides complement transfer without router discrimination; $\alpha = 0.1$ forces excessive routing.

\textbf{Model dimension $d$.} Tested at $d \in \{32, 64, 128, 256\}$. Results are qualitatively consistent; the group structure's advantage is proportionally similar across scales.

\textbf{Number range.} At very small num\_range ($\leq 10$), embeddings lack sufficient training context and all models underperform. At num\_range $\geq 15$, the sum function is learnable and the complement split provides clear signal.

\section{Scaling Analysis}

We conducted two scaling studies to determine whether the group advantage grows with problem size.

\textbf{num\_range sweep} (10--50, $d = 128$ and $d = 32$): all models converge to 100\% complement accuracy by num\_range $= 25$. The sum function decomposes into per-number values, so the effective table size grows as $O(n)$ rather than $O(n^3)$. This means num\_range scaling does not create the combinatorial challenge needed to separate the models.

\textbf{Coverage sweep} (1--4 orderings per triple, num\_range $= 15$): differentiation exists only at $k = 1$ (17\% coverage, 1:5 generalization). By $k = 2$ (33\%), all models converge to $>$99\%. The jump from $k = 1$ to $k = 2$ is the critical transition where even a single additional ordering provides enough context for memorization to succeed.

These results confirm that the sum function, while useful for validating the architecture, does not create a combinatorial scaling challenge. The right scaling axis is group order ($n!$ growth), not input range ($n$ growth).

\end{document}
